\label{sec:introduction}

In the 2011 California Citizens Redistricting Commission proceedings,
communities testified about their community bonds along six different
dimensions simultaneously.
Farmers in the San Joaquin Valley cited shared agricultural water rights
(economic), shared open farmland landscape (physical), shared agricultural
zoning districts (administrative), and shared markets (commuting shed).
Tech workers in Silicon Valley cited shared employment corridors (economic),
shared high-density commercial land use (physical), shared transit access
to Caltrain stations (transit), and shared municipal service districts (zone).
Coastal fishing communities in Monterey cited shared ocean access (physical),
shared fishing industry economy (economic), and shared county coastline
administration (zone).

Each of these communities was making a multi-dimensional argument.
The strength of the community bond was not that the community scored highest
on any single dimension but that it scored highly on \emph{multiple} dimensions
simultaneously---that the economic, physical, administrative, and behavioral
signals all agreed.

The M-track series, M.1 through M.7, defines seven individual community
character signals, each measured by a single cosine or Jaccard similarity
function on a different feature vector.
Each paper in the series provides a standalone signal that can be deployed
individually when the relevant data is available and the relevant community
type is present in the state being redistricted.

But the complete analytical tool is the composite index defined in this paper.
A plan drawn with all seven signals simultaneously produces communities that
are economically coherent, physically coherent, administratively coherent,
and transit-connected.
The composite index provides a single number that quantifies multi-dimensional
community cohesion for every adjacent tract pair in the state.

\paragraph{Why a composite rather than an ensemble.}
One alternative to the composite index is an ensemble approach: run
\textsc{bisect} with each individual signal separately and choose the plan
that satisfies the most criteria.
This is a reasonable approach but has two limitations:

\begin{enumerate}
  \item \textbf{Inconsistency.}
    A plan that is optimal for M.1 may be suboptimal for M.6.
    Selecting among plans that optimize different criteria requires
    a meta-criterion that is not clearly defined.
  \item \textbf{Legal exposure.}
    Selecting a plan because it scores well on a partisan criterion
    (even if that criterion is justified as a community-of-interest
    measure) creates legal exposure.
    The composite index approach selects a plan by simultaneously
    optimizing a weighted combination of non-partisan criteria,
    which is a cleaner legal position.
\end{enumerate}

The composite index resolves both limitations: it specifies the
meta-criterion (the weighted sum) explicitly and in advance of any
redistricting run, and it produces a single plan that can be defended
on all dimensions simultaneously.

\paragraph{The legal argument for a composite.}
Courts in redistricting cases have generally accepted evidence that a plan
preserves communities of interest when the evidence is specific, quantitative,
and non-partisan.
A composite index that is constructed from federal government data sources
(Census Bureau LODES, USGS NLCD, U.S.\ EIA Form 861, TIGER/Line) and
assigned weights that reflect the legal recognition of different community
types provides exactly this: specific (down to the tract pair), quantitative
(a number between 0 and 1), and non-partisan (no electoral data enters any
component).

\paragraph{Paper organization.}
Section~\ref{sec:background} defines the composite formula and its
normalization over available signals.
Section~\ref{sec:algorithm} justifies the default weights through a
calibration argument.
Section~\ref{sec:results} provides the complete court usage template,
including expert witness statement language for redistricting briefs
and testimony.
Section~\ref{sec:legal} describes the experimental design for Phase~2.
Section~\ref{sec:conclusion} concludes.
